\chapter{Headache}

\section*{Definition}
Pain perceived anywhere in head region among most common human experiences encountered clinically.

\section*{Pathophysiology}
Primary headaches lack identifiable structural cause; secondary headaches attributable to another condition with diverse mechanisms.

\section*{Causes}
\begin{itemize}
  \item Tension-type headache (most common)
  \item Migraine
  \item Cluster headache
  \item Medication-overuse
  \item Secondary causes requiring urgent evaluation
\end{itemize}

\section*{When to See a Doctor}
See your doctor if:
\begin{itemize}
  \item Severe or worsening symptoms
  \item Symptoms lasting more than Thunderclap onset, progressive worsening, new after age 50 require urgent evaluation
  \item Accompanied by other concerning signs
\end{itemize}

\section*{Related Symptoms}
\begin{itemize}
  \item Nausea/vomiting
  \item Photophobia/phonophobia
  \item Neck stiffness
  \item Visual aura changes
  \item Worsening with Valsalva
\end{itemize}
\textbf{Important:} If this symptom persists, your doctor will check for serious underlying causes.

\section*{Self-Care / Home Management}
\begin{itemize}
  \item Ensure a safe environment to prevent injury during episodes (remove hazards, avoid driving or operating machinery).
  \item Keep a symptom diary documenting triggers, duration, frequency, and associated features.
  \item Practice stress reduction techniques (deep breathing, meditation, adequate sleep) as stress often exacerbates neurological symptoms.
  \item Avoid alcohol, caffeine, and recreational drugs which can worsen many neurological conditions.
  \item Seek immediate evaluation for sudden-onset, severe, or progressive neurological symptoms.
\end{itemize}

\section*{How Your Doctor Will Evaluate You}
\begin{itemize}
  \item Detailed neurological history includes onset pattern, progression, triggers, associated symptoms, and functional impact.
  \item Comprehensive neurological examination assesses mental status, cranial nerves, motor/sensory function, reflexes, coordination, and gait.
  \item Neuroimaging (CT, MRI) is often indicated to evaluate for structural, vascular, or demyelinating causes.
  \item Specialized testing (EEG for seizures, nerve conduction studies for neuropathy, lumbar puncture for meningitis) is guided by clinical suspicion.
\end{itemize}

\section*{Treatment Options}
\begin{itemize}
  \item Treatment is directed at the underlying cause identified through diagnostic evaluation.
  \item Symptomatic pharmacotherapy may include analgesics, anticonvulsants, antidepressants, or disease-modifying agents as appropriate.
  \item Physical, occupational, and speech therapy play important roles in functional recovery and rehabilitation.
  \item Referral to neurology is indicated for unexplained, progressive, or treatment-refractory neurological symptoms.
\end{itemize}

\clearpage
